\begin{python0}
from solutions import *; clear()
\end{python0}
\sectionthree{Factorial}

First let's just do math.

Let f(n) denote the factorial function. In other words f(3) is 3x2x1;
f(5) is 5x4x3x2x1. In math the factorial of 3 is written 3!; the
factorial of 5 is written 5!. This is not new -- I have already talked
about this long time ago.

First of all there is a \EMPHASIZE{recursion} here. Look at f(n).

f(n) = n x (n-1) x (n-2) x ... x 3 x 2 x 1

Note that

f(n - 1) = (n-1) x (n-2) x ... x 3 x 2 x 1

Do you see f(n - 1) in f(n)? Putting the above together we have

f(n) = n x (n-1) x (n-2) x ... x 3 x 2 x 1

= n x [(n-1) x (n-2) x ... x 3 x 2 x 1]

in other words:

f(n) = n x f(n-1)

Right? You can use this mathematical fact to compute f(3) by hand:

f(3) = 3 x f(2)

f(2) = 2 x f(1)

WAIT!!! When do we stop?

f(1) = 1 x f(0)

f(0) = 0 x f(-1)

f(-1) = -1 x f(-2)

Not good ... this will go on forever. But hang on. Look at this:

f(3) = 3 x f(2)

f(2) = 2 x f(1)

f(1) = 1 x f(0)

This means

f(3) = 3 x f(2) = 3 x 2 x f(1) = 3 x 2 x 1 x f(0)

Now ... if we insist that f(0) is defined to be 1 we get

f(3) = 3 x 2 x 1 x f(0) = 3 x 2 x 1 x 1

So we have these properties for our factorial function:

f(n) = n x f(n-1), if n > 0

\EMPHASIZE{f(0) = 1}

The fact

\EMPHASIZE{f(0) = 1}

is called the \EMPHASIZE{base case} of f(n) while

f(n) = n x f(n-1), if n > 0

is called the \EMPHASIZE{recursive case}.

This is how you translate the above function:

f(n) = n x f(n-1), if n > 0

\EMPHASIZE{f(0) = 1}

into C++ code:

\begin{consolethree}[escapeinside=||]
int f(int n)
{
    if (n == 0)
        return 1;
    else
        return n * f(n - 1);
}

int main()
{
    for (int n = 0; n < 5; n++)
    {
        std::cout << n << "! = " << f(n)
                  << std::endl;
    }
    return 0;
}
\end{consolethree}

Of course there's another way (non-recursive way) of
computing factorials \ldots{} using loops:

\begin{consolethree}[escapeinside=||]
int f(int n)
{
    int p = 1;
    for (int i = 1; i <= n; i++)
    {
        p *= i;
    }
    return p;
}
\end{consolethree}

Note that just like the sum(n) example the \EMPHASIZE{recursive version
does not have loops}.

To see the program trace its steps run this:

\begin{consolethree}[escapeinside=||]
int f(int n)
{
    std::cout << "enter f(" << n << ")" << std::endl;
    if (n == 0)
    {
        std::cout << "base case ... returning 1"
                  << std::endl;
        return 1;
    }
    else
    {
        int x = n * f(n - 1);
        std::cout << "recursive case ... returning "
                  << x << std::endl;
        return x;
    }
}

int main()
{
    int x = f(5);
    std::cout << "main() ... " << x << std::endl;
    return 0;
}
\end{consolethree}

Here's a diagram that might help:
\begin{python}
from latextool_basic import *
p = Plot()

p += Rect(0, 8, 3, 9, label=r'\texttt{f(0)}', linewidth=0.05)
p += Rect(0, 6, 3, 7, label=r'\texttt{f(1)}', linewidth=0.05)
p += Rect(0, 4, 3, 5, label=r'\texttt{f(2)}', linewidth=0.05)
p += Rect(0, 2, 3, 3, label=r'\texttt{f(3)}', linewidth=0.05)
p += Rect(0, 0, 3, 1, label=r'\texttt{f(4)}', linewidth=0.05)
#up
zaphod = -1
for i in range(5):
    zaphod += 1
    p += Line(points=[(0.75, zaphod - 1), (0.75, zaphod)], linewidth=0.075, endstyle='>')
    zaphod += 1

namez = ['1', '1', '2', '6', '24']
#down
for i in range(5):
    zaphod -= 1
    p += Line(points=[(2.25, zaphod), (2.25, zaphod - 1)], linewidth=0.075, endstyle='>', label= r'\tab[5em]{' + namez[i] + '}')
    zaphod -= 1

p += Rect(5, -0.5, 14, 7.5, linewidth=0.1, innersep=0.2, s=r"""In \texttt{f(3)}, \begin{console}
int f(int n)
{
    if (n == 0)
       return 1;
    else
       return n * f(n - 1); 
}
\end{console}
On receiving \texttt{f(2)}, i.e. 2, \texttt{f(3)} returns \texttt{n * f(n - 1)}, i.e. 3 * 2 i.e. 6.""", align='c')
p += Line(points=[(2.5, 3), (5, 3)], linestyle='dashed', linewidth=0.1)
print(p)
\end{python}


Notice that for the function call graph of f(4), there are 5 boxes. You
should see quickly that the function call graph for f(n) has n + 1
boxes.


